\documentclass[12pt, a4paper]{ctexart}
\usepackage[margin=2cm]{geometry}
\usepackage{enumitem}
\usepackage{fancyhdr}
\usepackage{amsmath, amssymb}
\usepackage{graphicx}
\usepackage{xcolor}

\setCJKmainfont{Hiragino Sans GB}
\setCJKsansfont{Heiti SC}

\pagestyle{fancy}
\fancyhf{}
\fancyhead[L]{八年级下册数学测试卷}
\fancyhead[R]{第~\thepage~页}
\fancyfoot[C]{}
\renewcommand{\headrulewidth}{0.4pt}

\begin{document}

\begin{center}
    {\LARGE\bfseries 八年级下册数学测试卷}\\[0.5cm]
    {\large 考试时间：120分钟 \quad 满分：120分}\\[0.3cm]
    {\normalsize 注意：请将答案写在答题卡上，写在试卷上无效。}
\end{center}

\vspace{0.5cm}

{\large\bfseries 一、选择题（本大题共10小题，每小题3分，共30分。每小题只有一个正确选项。）}

\begin{enumerate}[label=\arabic*.]

\item $\sqrt{9}$的值是（\hspace{1cm}）
\begin{enumerate}[label=\Alph*.]
\quad \textbf{A.}~$3$ \quad \textbf{B.}~$-3$ \quad \textbf{C.}~$\pm 3$ \quad \textbf{D.}~$81$
\end{enumerate}

\item 下列各组数中，能构成直角三角形三边的是（\hspace{1cm}）
\begin{enumerate}[label=\Alph*.]
\quad \textbf{A.}~$1,2,3$ \quad \textbf{B.}~$3,4,6$ \quad \textbf{C.}~$5,12,13$ \quad \textbf{D.}~$6,7,10$
\end{enumerate}

\item 下列二次根式中，是最简二次根式的是（\hspace{1cm}）
\begin{enumerate}[label=\Alph*.]
\quad \textbf{A.}~$\sqrt{8}$ \quad \textbf{B.}~$\sqrt{12}$ \quad \textbf{C.}~$\sqrt{15}$ \quad \textbf{D.}~$\sqrt{\dfrac{1}{3}}$
\end{enumerate}

\item 平行四边形$ABCD$中，$\angle A=110^{\circ}$，则$\angle B$的度数为（\hspace{1cm}）
\begin{enumerate}[label=\Alph*.]
\quad \textbf{A.}~$70^{\circ}$ \quad \textbf{B.}~$110^{\circ}$ \quad \textbf{C.}~$60^{\circ}$ \quad \textbf{D.}~$120^{\circ}$
\end{enumerate}

\item 下列函数中，是一次函数的是（\hspace{1cm}）
\begin{enumerate}[label=\Alph*.]
\quad \textbf{A.}~$y=\dfrac{1}{x}$ \quad \textbf{B.}~$y=3x^{2}+1$ \quad \textbf{C.}~$y=2x-5$ \quad \textbf{D.}~$y=\sqrt{x}$
\end{enumerate}

\item 一组数据$2,3,5,5,7,8$的中位数和众数分别是（\hspace{1cm}）
\begin{enumerate}[label=\Alph*.]
\quad \textbf{A.}~$5,5$ \quad \textbf{B.}~$5,3$ \quad \textbf{C.}~$4,5$ \quad \textbf{D.}~$5,4$
\end{enumerate}

\item 化简$\sqrt{12}-\sqrt{3}$的结果是（\hspace{1cm}）
\begin{enumerate}[label=\Alph*.]
\quad \textbf{A.}~$\sqrt{3}$ \quad \textbf{B.}~$3\sqrt{3}$ \quad \textbf{C.}~$2\sqrt{3}$ \quad \textbf{D.}~$3$
\end{enumerate}

\item 正方形的对角线长为$4\sqrt{2}$，则正方形的边长为（\hspace{1cm}）
\begin{enumerate}[label=\Alph*.]
\quad \textbf{A.}~$4$ \quad \textbf{B.}~$4\sqrt{2}$ \quad \textbf{C.}~$8$ \quad \textbf{D.}~$2\sqrt{2}$
\end{enumerate}

\item 一次函数$y=2x-3$与$x$轴的交点坐标为（\hspace{1cm}）
\begin{enumerate}[label=\Alph*.]
\quad \textbf{A.}~$(0,-3)$ \quad \textbf{B.}~$\left(\dfrac{3}{2},0\right)$ \quad \textbf{C.}~$(3,0)$ \quad \textbf{D.}~$\left(0,\dfrac{3}{2}\right)$
\end{enumerate}

\item 若一组数据$x_{1},x_{2},x_{3},x_{4},x_{5}$的平均数为$8$，方差为$2$，则数据$2x_{1}+1,2x_{2}+1,2x_{3}+1,2x_{4}+1,2x_{5}+1$的平均数和方差分别为（\hspace{1cm}）
\begin{enumerate}[label=\Alph*.]
\quad \textbf{A.}~$17,4$ \quad \textbf{B.}~$17,8$ \quad \textbf{C.}~$16,4$ \quad \textbf{D.}~$16,8$
\end{enumerate}

\end{enumerate}

\vspace{0.5cm}

{\large\bfseries 二、填空题（本大题共6小题，每小题3分，共18分。）}

\begin{enumerate}[label=\arabic*., start=11]

\item 化简：$\sqrt{18}=$ \underline{\hspace{3cm}} 。

\item 在$\mathrm{Rt}\triangle ABC$中，$\angle C=90^{\circ}$，$AC=5$，$BC=12$，则$AB=$ \underline{\hspace{3cm}} 。

\item 直线$y=2x+1$与$y$轴的交点坐标为 \underline{\hspace{3cm}} 。

\item 如图，平行四边形$ABCD$的对角线$AC$和$BD$相交于点$O$，若$AC=6$，$BD=8$，$AC\perp BD$，则$AB=$ \underline{\hspace{3cm}} 。

\item 已知$y-3$与$x$成正比例关系，且当$x=2$时，$y=7$，则$y$与$x$的函数关系式为 \underline{\hspace{3cm}} 。

\item 数据$1,2,3,x,5$的平均数为$3$，则$x=$ \underline{\hspace{3cm}} ，这组数据的方差为 \underline{\hspace{3cm}} 。

\end{enumerate}

\vspace{0.5cm}

{\large\bfseries 三、解答题（本大题共8小题，共72分。解答应写出文字说明、演算步骤或证明过程。）}

\begin{enumerate}[label=\arabic*., start=17]

\item （8分）计算：

（1）$\sqrt{48}-\sqrt{27}+\sqrt{3}$；

（2）$(2\sqrt{3}-\sqrt{2})(2\sqrt{3}+\sqrt{2})$。

\vspace{6cm}

\item （8分）如图，在$\triangle ABC$中，$\angle C=90^{\circ}$，$AC=8$，$BC=15$。

（1）求$AB$的长；

（2）求$AB$边上的高$CD$的长。

\vspace{6cm}

\item （8分）如图，在平行四边形$ABCD$中，$E$、$F$分别是$AD$、$BC$的中点。

求证：$BE=DF$。

\vspace{6cm}

\item （8分）已知一次函数$y=kx+b$的图象经过点$A(0,-3)$和点$B(2,1)$。

（1）求该一次函数的解析式；

（2）求该函数图象与$x$轴的交点坐标。

\vspace{6cm}

\item （10分）某校对八年级学生进行体育测试，$20$名同学的跳绳成绩（单位：次/分钟）如下：

$120,135,140,145,150,150,155,160,165,170,$

$170,170,175,180,180,185,190,195,200,210$

（1）求这组数据的平均数；

（2）求这组数据的中位数和众数。

\vspace{6cm}

\item （10分）如图，在$\triangle ABC$中，$\angle C=90^{\circ}$，$AC=6$，$BC=8$，$CD\perp AB$于$D$。

（1）求$AB$的长；

（2）求$CD$的长。

\vspace{6cm}

\item （10分）如图，矩形$ABCD$中，$AB=6$，$BC=8$，将矩形沿对角线$BD$折叠，点$C$落在点$E$处，$BE$交$AD$于点$F$。

（1）求$BD$的长；

（2）求$BF$的长。

\vspace{6cm}

\item （10分）某工厂有甲、乙两个车间，甲车间每天生产$A$产品$60$件，$B$产品$40$件；乙车间每天生产$A$产品$30$件，$B$产品$50$件。现需生产$A$产品$300$件，$B$产品$340$件。

（1）设甲车间工作$x$天，乙车间工作$y$天，列出关于$x$、$y$的方程组；

（2）求甲、乙两车间各需工作多少天。

\vspace{6cm}

\end{enumerate}

\end{document}
